\usepackage[T1]{fontenc}
\usepackage[utf8]{inputenc}
\usepackage{geometry}
\geometry{verbose,tmargin=1in,bmargin=1in,lmargin=1in,rmargin=1in}
\usepackage{float}
\usepackage{booktabs}
\usepackage{tabularx}
\usepackage{amsmath}
\usepackage{amsthm}
\usepackage{amssymb}
\usepackage{bm}
\usepackage{graphicx}
\usepackage{units}
\usepackage{enumitem}
\usepackage{microtype}          % typographic polish: char protrusion, expansion
\usepackage{xcolor}
\usepackage{comment}
\usepackage{tcolorbox}          % highlighted result boxes
\tcbuselibrary{skins,breakable}
\usepackage{fancyhdr}           % running header with the short title
\usepackage[textsize=footnotesize]{todonotes}   % margin \todo{} notes; width follows \marginparwidth
\usepackage{authblk}

\makeatletter

\numberwithin{equation}{section}
\numberwithin{figure}{section}
\numberwithin{table}{section}

% ---- Theorem environments ----
\@ifundefined{thm}{%
  \theoremstyle{plain}%
  \newtheorem{thm}{Theorem}[section]%
}{}
\@ifundefined{prop}{\newtheorem{prop}[thm]{Proposition}}{}
\@ifundefined{lem}{\newtheorem{lem}[thm]{Lemma}}{}
\@ifundefined{cor}{\newtheorem{cor}[thm]{Corollary}}{}
\@ifundefined{conj}{\newtheorem{conj}[thm]{Conjecture}}{}
\@ifundefined{defn}{%
  \theoremstyle{definition}%
  \newtheorem{defn}[thm]{Definition}%
}{}
\@ifundefined{assum}{%
  \theoremstyle{definition}%
  \newtheorem{assum}[thm]{Assumption}%
}{}
\@ifundefined{example}{\newtheorem{example}[thm]{Example}}{}
\@ifundefined{remark}{%
  \theoremstyle{remark}%
  \newtheorem{remark}[thm]{Remark}%
}{}

% ---- Convenient macros ----
\newcommand{\R}{\mathbb{R}}
\newcommand{\C}{\mathbb{C}}
\newcommand{\N}{\mathbb{N}}
\newcommand{\E}{\mathbb{E}}
\newcommand{\Prob}{\mathbb{P}}
\newcommand{\ip}[2]{\langle #1, #2 \rangle}
\newcommand{\norm}[1]{\left\| #1 \right\|}
\newcommand{\abs}[1]{\left| #1 \right|}
\newcommand{\T}{^{\mathsf{T}}}          % transpose
\newcommand{\dd}{\,\mathrm{d}}          % upright differential
\newcommand{\D}[2]{\frac{\mathrm{d}#1}{\mathrm{d}#2}}     % total derivative
\newcommand{\pd}[2]{\frac{\partial #1}{\partial #2}}      % partial derivative

% ---- Program-wide notation ----
\newcommand{\Rset}{\mathcal{R}}          % reachable set
\newcommand{\Mplus}{\mathcal{M}_+}       % cone of nonnegative Radon measures
\newcommand{\Xset}{\mathcal{X}}          % state constraint set
\newcommand{\Uset}{\mathcal{U}}          % control constraint set
\newcommand{\Lie}{\mathcal{L}}           % Lie derivative / generator
\newcommand{\modes}{\mathcal{Q}}         % mode alphabet (Paper 1, Section 2.2)
\newcommand{\guard}{\mathcal{S}}         % guard surface
\newcommand{\reset}{\mathcal{J}}         % reset map

% The Henrion-Korda dual pair is conventionally (v, w). Both letters are
% already body-frame velocity components in the state vector, and the
% certificate is a function OF that state, so the two would collide inside a
% single displayed equation. Capital Greek instead: distinct from latitude
% \phi and heading \psi, which are lowercase. Named here so the choice is
% changed in one place rather than hunted through the papers.
\newcommand{\cert}{\Phi}               % Liouville supersolution, cert(t,x)
\newcommand{\certT}{\Psi}              % terminal majorant, certT(x)

\makeatother

% ---- Highlighted-result box ----
\newtcolorbox{keybox}[1][]{colback=blue!4,colframe=blue!45!black,
  boxrule=0.6pt,arc=2pt,left=6pt,right=6pt,top=4pt,bottom=4pt,#1}

% ---- Development annotations ----
%  Three coloured boxes carrying the drafting ledger inside the source:
%    decide -- a choice made, recorded so it is not silently re-opened
%    gap    -- an obligation not yet discharged
%    idea   -- a line of argument worth keeping, not yet load-bearing
%  They are deliberately conspicuous. Set \draftnotesfalse below (or pass
%  \PassOptionsToPackage before \usepackage[T1]{fontenc}
\usepackage[utf8]{inputenc}
\usepackage{geometry}
\geometry{verbose,tmargin=1in,bmargin=1in,lmargin=1in,rmargin=1in}
\usepackage{float}
\usepackage{booktabs}
\usepackage{tabularx}
\usepackage{amsmath}
\usepackage{amsthm}
\usepackage{amssymb}
\usepackage{bm}
\usepackage{graphicx}
\usepackage{units}
\usepackage{enumitem}
\usepackage{microtype}          % typographic polish: char protrusion, expansion
\usepackage{xcolor}
\usepackage{comment}
\usepackage{tcolorbox}          % highlighted result boxes
\tcbuselibrary{skins,breakable}
\usepackage{fancyhdr}           % running header with the short title
\usepackage[textsize=footnotesize]{todonotes}   % margin \todo{} notes; width follows \marginparwidth
\usepackage{authblk}

\makeatletter

\numberwithin{equation}{section}
\numberwithin{figure}{section}
\numberwithin{table}{section}

% ---- Theorem environments ----
\@ifundefined{thm}{%
  \theoremstyle{plain}%
  \newtheorem{thm}{Theorem}[section]%
}{}
\@ifundefined{prop}{\newtheorem{prop}[thm]{Proposition}}{}
\@ifundefined{lem}{\newtheorem{lem}[thm]{Lemma}}{}
\@ifundefined{cor}{\newtheorem{cor}[thm]{Corollary}}{}
\@ifundefined{conj}{\newtheorem{conj}[thm]{Conjecture}}{}
\@ifundefined{defn}{%
  \theoremstyle{definition}%
  \newtheorem{defn}[thm]{Definition}%
}{}
\@ifundefined{assum}{%
  \theoremstyle{definition}%
  \newtheorem{assum}[thm]{Assumption}%
}{}
\@ifundefined{example}{\newtheorem{example}[thm]{Example}}{}
\@ifundefined{remark}{%
  \theoremstyle{remark}%
  \newtheorem{remark}[thm]{Remark}%
}{}

% ---- Convenient macros ----
\newcommand{\R}{\mathbb{R}}
\newcommand{\C}{\mathbb{C}}
\newcommand{\N}{\mathbb{N}}
\newcommand{\E}{\mathbb{E}}
\newcommand{\Prob}{\mathbb{P}}
\newcommand{\ip}[2]{\langle #1, #2 \rangle}
\newcommand{\norm}[1]{\left\| #1 \right\|}
\newcommand{\abs}[1]{\left| #1 \right|}
\newcommand{\T}{^{\mathsf{T}}}          % transpose
\newcommand{\dd}{\,\mathrm{d}}          % upright differential
\newcommand{\D}[2]{\frac{\mathrm{d}#1}{\mathrm{d}#2}}     % total derivative
\newcommand{\pd}[2]{\frac{\partial #1}{\partial #2}}      % partial derivative

% ---- Program-wide notation ----
\newcommand{\Rset}{\mathcal{R}}          % reachable set
\newcommand{\Mplus}{\mathcal{M}_+}       % cone of nonnegative Radon measures
\newcommand{\Xset}{\mathcal{X}}          % state constraint set
\newcommand{\Uset}{\mathcal{U}}          % control constraint set
\newcommand{\Lie}{\mathcal{L}}           % Lie derivative / generator
\newcommand{\modes}{\mathcal{Q}}         % mode alphabet (Paper 1, Section 2.2)
\newcommand{\guard}{\mathcal{S}}         % guard surface
\newcommand{\reset}{\mathcal{J}}         % reset map

% The Henrion-Korda dual pair is conventionally (v, w). Both letters are
% already body-frame velocity components in the state vector, and the
% certificate is a function OF that state, so the two would collide inside a
% single displayed equation. Capital Greek instead: distinct from latitude
% \phi and heading \psi, which are lowercase. Named here so the choice is
% changed in one place rather than hunted through the papers.
\newcommand{\cert}{\Phi}               % Liouville supersolution, cert(t,x)
\newcommand{\certT}{\Psi}              % terminal majorant, certT(x)

\makeatother

% ---- Highlighted-result box ----
\newtcolorbox{keybox}[1][]{colback=blue!4,colframe=blue!45!black,
  boxrule=0.6pt,arc=2pt,left=6pt,right=6pt,top=4pt,bottom=4pt,#1}

% ---- Development annotations ----
%  Three coloured boxes carrying the drafting ledger inside the source:
%    decide -- a choice made, recorded so it is not silently re-opened
%    gap    -- an obligation not yet discharged
%    idea   -- a line of argument worth keeping, not yet load-bearing
%  They are deliberately conspicuous. Set \draftnotesfalse below (or pass
%  \PassOptionsToPackage before \usepackage[T1]{fontenc}
\usepackage[utf8]{inputenc}
\usepackage{geometry}
\geometry{verbose,tmargin=1in,bmargin=1in,lmargin=1in,rmargin=1in}
\usepackage{float}
\usepackage{booktabs}
\usepackage{tabularx}
\usepackage{amsmath}
\usepackage{amsthm}
\usepackage{amssymb}
\usepackage{bm}
\usepackage{graphicx}
\usepackage{units}
\usepackage{enumitem}
\usepackage{microtype}          % typographic polish: char protrusion, expansion
\usepackage{xcolor}
\usepackage{comment}
\usepackage{tcolorbox}          % highlighted result boxes
\tcbuselibrary{skins,breakable}
\usepackage{fancyhdr}           % running header with the short title
\usepackage[textsize=footnotesize]{todonotes}   % margin \todo{} notes; width follows \marginparwidth
\usepackage{authblk}

\makeatletter

\numberwithin{equation}{section}
\numberwithin{figure}{section}
\numberwithin{table}{section}

% ---- Theorem environments ----
\@ifundefined{thm}{%
  \theoremstyle{plain}%
  \newtheorem{thm}{Theorem}[section]%
}{}
\@ifundefined{prop}{\newtheorem{prop}[thm]{Proposition}}{}
\@ifundefined{lem}{\newtheorem{lem}[thm]{Lemma}}{}
\@ifundefined{cor}{\newtheorem{cor}[thm]{Corollary}}{}
\@ifundefined{conj}{\newtheorem{conj}[thm]{Conjecture}}{}
\@ifundefined{defn}{%
  \theoremstyle{definition}%
  \newtheorem{defn}[thm]{Definition}%
}{}
\@ifundefined{assum}{%
  \theoremstyle{definition}%
  \newtheorem{assum}[thm]{Assumption}%
}{}
\@ifundefined{example}{\newtheorem{example}[thm]{Example}}{}
\@ifundefined{remark}{%
  \theoremstyle{remark}%
  \newtheorem{remark}[thm]{Remark}%
}{}

% ---- Convenient macros ----
\newcommand{\R}{\mathbb{R}}
\newcommand{\C}{\mathbb{C}}
\newcommand{\N}{\mathbb{N}}
\newcommand{\E}{\mathbb{E}}
\newcommand{\Prob}{\mathbb{P}}
\newcommand{\ip}[2]{\langle #1, #2 \rangle}
\newcommand{\norm}[1]{\left\| #1 \right\|}
\newcommand{\abs}[1]{\left| #1 \right|}
\newcommand{\T}{^{\mathsf{T}}}          % transpose
\newcommand{\dd}{\,\mathrm{d}}          % upright differential
\newcommand{\D}[2]{\frac{\mathrm{d}#1}{\mathrm{d}#2}}     % total derivative
\newcommand{\pd}[2]{\frac{\partial #1}{\partial #2}}      % partial derivative

% ---- Program-wide notation ----
\newcommand{\Rset}{\mathcal{R}}          % reachable set
\newcommand{\Mplus}{\mathcal{M}_+}       % cone of nonnegative Radon measures
\newcommand{\Xset}{\mathcal{X}}          % state constraint set
\newcommand{\Uset}{\mathcal{U}}          % control constraint set
\newcommand{\Lie}{\mathcal{L}}           % Lie derivative / generator
\newcommand{\modes}{\mathcal{Q}}         % mode alphabet (Paper 1, Section 2.2)
\newcommand{\guard}{\mathcal{S}}         % guard surface
\newcommand{\reset}{\mathcal{J}}         % reset map

% The Henrion-Korda dual pair is conventionally (v, w). Both letters are
% already body-frame velocity components in the state vector, and the
% certificate is a function OF that state, so the two would collide inside a
% single displayed equation. Capital Greek instead: distinct from latitude
% \phi and heading \psi, which are lowercase. Named here so the choice is
% changed in one place rather than hunted through the papers.
\newcommand{\cert}{\Phi}               % Liouville supersolution, cert(t,x)
\newcommand{\certT}{\Psi}              % terminal majorant, certT(x)

\makeatother

% ---- Highlighted-result box ----
\newtcolorbox{keybox}[1][]{colback=blue!4,colframe=blue!45!black,
  boxrule=0.6pt,arc=2pt,left=6pt,right=6pt,top=4pt,bottom=4pt,#1}

% ---- Development annotations ----
%  Three coloured boxes carrying the drafting ledger inside the source:
%    decide -- a choice made, recorded so it is not silently re-opened
%    gap    -- an obligation not yet discharged
%    idea   -- a line of argument worth keeping, not yet load-bearing
%  They are deliberately conspicuous. Set \draftnotesfalse below (or pass
%  \PassOptionsToPackage before \usepackage[T1]{fontenc}
\usepackage[utf8]{inputenc}
\usepackage{geometry}
\geometry{verbose,tmargin=1in,bmargin=1in,lmargin=1in,rmargin=1in}
\usepackage{float}
\usepackage{booktabs}
\usepackage{tabularx}
\usepackage{amsmath}
\usepackage{amsthm}
\usepackage{amssymb}
\usepackage{bm}
\usepackage{graphicx}
\usepackage{units}
\usepackage{enumitem}
\usepackage{microtype}          % typographic polish: char protrusion, expansion
\usepackage{xcolor}
\usepackage{comment}
\usepackage{tcolorbox}          % highlighted result boxes
\tcbuselibrary{skins,breakable}
\usepackage{fancyhdr}           % running header with the short title
\usepackage[textsize=footnotesize]{todonotes}   % margin \todo{} notes; width follows \marginparwidth
\usepackage{authblk}

\makeatletter

\numberwithin{equation}{section}
\numberwithin{figure}{section}
\numberwithin{table}{section}

% ---- Theorem environments ----
\@ifundefined{thm}{%
  \theoremstyle{plain}%
  \newtheorem{thm}{Theorem}[section]%
}{}
\@ifundefined{prop}{\newtheorem{prop}[thm]{Proposition}}{}
\@ifundefined{lem}{\newtheorem{lem}[thm]{Lemma}}{}
\@ifundefined{cor}{\newtheorem{cor}[thm]{Corollary}}{}
\@ifundefined{conj}{\newtheorem{conj}[thm]{Conjecture}}{}
\@ifundefined{defn}{%
  \theoremstyle{definition}%
  \newtheorem{defn}[thm]{Definition}%
}{}
\@ifundefined{assum}{%
  \theoremstyle{definition}%
  \newtheorem{assum}[thm]{Assumption}%
}{}
\@ifundefined{example}{\newtheorem{example}[thm]{Example}}{}
\@ifundefined{remark}{%
  \theoremstyle{remark}%
  \newtheorem{remark}[thm]{Remark}%
}{}

% ---- Convenient macros ----
\newcommand{\R}{\mathbb{R}}
\newcommand{\C}{\mathbb{C}}
\newcommand{\N}{\mathbb{N}}
\newcommand{\E}{\mathbb{E}}
\newcommand{\Prob}{\mathbb{P}}
\newcommand{\ip}[2]{\langle #1, #2 \rangle}
\newcommand{\norm}[1]{\left\| #1 \right\|}
\newcommand{\abs}[1]{\left| #1 \right|}
\newcommand{\T}{^{\mathsf{T}}}          % transpose
\newcommand{\dd}{\,\mathrm{d}}          % upright differential
\newcommand{\D}[2]{\frac{\mathrm{d}#1}{\mathrm{d}#2}}     % total derivative
\newcommand{\pd}[2]{\frac{\partial #1}{\partial #2}}      % partial derivative

% ---- Program-wide notation ----
\newcommand{\Rset}{\mathcal{R}}          % reachable set
\newcommand{\Mplus}{\mathcal{M}_+}       % cone of nonnegative Radon measures
\newcommand{\Xset}{\mathcal{X}}          % state constraint set
\newcommand{\Uset}{\mathcal{U}}          % control constraint set
\newcommand{\Lie}{\mathcal{L}}           % Lie derivative / generator
\newcommand{\modes}{\mathcal{Q}}         % mode alphabet (Paper 1, Section 2.2)
\newcommand{\guard}{\mathcal{S}}         % guard surface
\newcommand{\reset}{\mathcal{J}}         % reset map

% The Henrion-Korda dual pair is conventionally (v, w). Both letters are
% already body-frame velocity components in the state vector, and the
% certificate is a function OF that state, so the two would collide inside a
% single displayed equation. Capital Greek instead: distinct from latitude
% \phi and heading \psi, which are lowercase. Named here so the choice is
% changed in one place rather than hunted through the papers.
\newcommand{\cert}{\Phi}               % Liouville supersolution, cert(t,x)
\newcommand{\certT}{\Psi}              % terminal majorant, certT(x)

\makeatother

% ---- Highlighted-result box ----
\newtcolorbox{keybox}[1][]{colback=blue!4,colframe=blue!45!black,
  boxrule=0.6pt,arc=2pt,left=6pt,right=6pt,top=4pt,bottom=4pt,#1}

% ---- Development annotations ----
%  Three coloured boxes carrying the drafting ledger inside the source:
%    decide -- a choice made, recorded so it is not silently re-opened
%    gap    -- an obligation not yet discharged
%    idea   -- a line of argument worth keeping, not yet load-bearing
%  They are deliberately conspicuous. Set \draftnotesfalse below (or pass
%  \PassOptionsToPackage before \input{preamble}) to compile a submission
%  copy in which all three vanish, exactly as [disable]{todonotes} removes
%  the \todo list --- the boxes must not survive into a submitted draft,
%  and a single switch is the only reliable way to be sure of that.
\newif\ifdraftnotes
\draftnotestrue

\newtcolorbox{decide@box}[1][]{colback=orange!4,colframe=orange!45!black,
  boxrule=0.6pt,arc=2pt,left=6pt,right=6pt,top=4pt,bottom=4pt,#1}
\newtcolorbox{gap@box}[1][]{colback=red!4,colframe=red!50!black,
  boxrule=0.6pt,arc=2pt,left=6pt,right=6pt,top=4pt,bottom=4pt,#1}
\newtcolorbox{idea@box}[1][]{colback=green!4,colframe=green!35!black,
  boxrule=0.6pt,arc=2pt,left=6pt,right=6pt,top=4pt,bottom=4pt,#1}

\ifdraftnotes
  \newenvironment{decide}{\begin{decide@box}}{\end{decide@box}}
  \newenvironment{gap}{\begin{gap@box}}{\end{gap@box}}
  \newenvironment{idea}{\begin{idea@box}}{\end{idea@box}}
\else
  \excludecomment{decide}
  \excludecomment{gap}
  \excludecomment{idea}
\fi

% ---- Theorem-environment aliases ----
%  The drafts write \begin{assumption} and \begin{lemma} in full; the
%  counters above are declared as assum and lem. Alias rather than rename,
%  so neither spelling is a compile error and no draft has to be rewritten
%  to match a decision about abbreviation.
\makeatletter
\@ifundefined{assumption}{%
  \newenvironment{assumption}[1][]{\begin{assum}[#1]}{\end{assum}}}{}
\@ifundefined{lemma}{%
  \newenvironment{lemma}[1][]{\begin{lem}[#1]}{\end{lem}}}{}
\makeatother

\usepackage{babel}

% ---- Cross-references and links (load late: hyperref then cleveref) ----
\usepackage[hidelinks,colorlinks=true,
  linkcolor=blue!55!black,citecolor=blue!55!black,urlcolor=blue!55!black]{hyperref}

% ---- Bibliography ----
%  biblatex, not the plain BibTeX \bibliography. Loaded after hyperref and
%  before cleveref, which is the order the biblatex manual asks for.
%
%  style=numeric keeps the bare-number in-text appearance the papers have
%  always used. sorting=none orders the reference list by first citation,
%  which reads better than alphabetical in a proposal.
%
%  sortcites is the one non-obvious option. It sorts the keys WITHIN each
%  citation so the numeric style can collapse consecutive labels into a
%  range -- without it a five-key \parencite renders as five separate
%  brackets. Note this is NOT natbib's sort&compress, which biblatex does
%  not provide: the biblatex natbib compatibility layer (blx-natbib.def)
%  defines \citep and \citet but nothing for compression, and an unknown
%  package key only warns, so a wrong setting here fails silently.
%
%  maxcitenames=2 keeps a multi-key citation from printing every author;
%  maxbibnames=99 with giveninits stops the reference list collapsing long
%  author lists to "et al." in a way that would not match the full lists
%  written out in shared.bib (see emmert2021nrlmsis, hersbach2020era5).
%
%  The per-entry reading notes in the bibliography live in the biblatex
%  `annotation' field. They are the most valuable content in that file and are
%  deliberately NOT printed. Note that `annotation' is a DATA field, not an
%  option of \printbibliography -- passing annotation=false to \printbibliography
%  is an xkeyval error, not a way of suppressing it. Blanking the field's format
%  is the mechanism, and it is the right one: the notes survive in the .bib and
%  in the .bbl for the author's use, and never reach the page. This must come
%  AFTER the \usepackage line below, since \DeclareFieldFormat is defined there.
\usepackage[backend=biber,style=numeric,sorting=none,sortcites,
  maxbibnames=99,giveninits=true,maxcitenames=2]{biblatex}

% biblatex requires csquotes whenever babel is loaded, and warns without it
% ("'babel/polyglossia' detected but 'csquotes' missing"). It supplies the
% quote punctuation the bibliography macros expect. Not currently used by any
% entry, but the dependency is unconditional.
\usepackage{csquotes}

\DeclareFieldFormat{annotation}{}

% The thesis bibliography. \addbibresource must be issued in the preamble, and
% this preamble is \input from a directory one level below manuscript/, so the
% path is written relative to that directory, not to this file.
\addbibresource{../thesis.bib}

\usepackage[nameinlink,capitalise]{cleveref}

% ---- Running header: each paper calls \runninghead{Short Title} ----
\newcommand{\runninghead}[1]{%
  \pagestyle{fancy}%
  \fancyhf{}%
  \lhead{\small #1}%
  \rhead{\small\thepage}%
  \renewcommand{\headrulewidth}{0.3pt}%
}
) to compile a submission
%  copy in which all three vanish, exactly as [disable]{todonotes} removes
%  the \todo list --- the boxes must not survive into a submitted draft,
%  and a single switch is the only reliable way to be sure of that.
\newif\ifdraftnotes
\draftnotestrue

\newtcolorbox{decide@box}[1][]{colback=orange!4,colframe=orange!45!black,
  boxrule=0.6pt,arc=2pt,left=6pt,right=6pt,top=4pt,bottom=4pt,#1}
\newtcolorbox{gap@box}[1][]{colback=red!4,colframe=red!50!black,
  boxrule=0.6pt,arc=2pt,left=6pt,right=6pt,top=4pt,bottom=4pt,#1}
\newtcolorbox{idea@box}[1][]{colback=green!4,colframe=green!35!black,
  boxrule=0.6pt,arc=2pt,left=6pt,right=6pt,top=4pt,bottom=4pt,#1}

\ifdraftnotes
  \newenvironment{decide}{\begin{decide@box}}{\end{decide@box}}
  \newenvironment{gap}{\begin{gap@box}}{\end{gap@box}}
  \newenvironment{idea}{\begin{idea@box}}{\end{idea@box}}
\else
  \excludecomment{decide}
  \excludecomment{gap}
  \excludecomment{idea}
\fi

% ---- Theorem-environment aliases ----
%  The drafts write \begin{assumption} and \begin{lemma} in full; the
%  counters above are declared as assum and lem. Alias rather than rename,
%  so neither spelling is a compile error and no draft has to be rewritten
%  to match a decision about abbreviation.
\makeatletter
\@ifundefined{assumption}{%
  \newenvironment{assumption}[1][]{\begin{assum}[#1]}{\end{assum}}}{}
\@ifundefined{lemma}{%
  \newenvironment{lemma}[1][]{\begin{lem}[#1]}{\end{lem}}}{}
\makeatother

\usepackage{babel}

% ---- Cross-references and links (load late: hyperref then cleveref) ----
\usepackage[hidelinks,colorlinks=true,
  linkcolor=blue!55!black,citecolor=blue!55!black,urlcolor=blue!55!black]{hyperref}

% ---- Bibliography ----
%  biblatex, not the plain BibTeX \bibliography. Loaded after hyperref and
%  before cleveref, which is the order the biblatex manual asks for.
%
%  style=numeric keeps the bare-number in-text appearance the papers have
%  always used. sorting=none orders the reference list by first citation,
%  which reads better than alphabetical in a proposal.
%
%  sortcites is the one non-obvious option. It sorts the keys WITHIN each
%  citation so the numeric style can collapse consecutive labels into a
%  range -- without it a five-key \parencite renders as five separate
%  brackets. Note this is NOT natbib's sort&compress, which biblatex does
%  not provide: the biblatex natbib compatibility layer (blx-natbib.def)
%  defines \citep and \citet but nothing for compression, and an unknown
%  package key only warns, so a wrong setting here fails silently.
%
%  maxcitenames=2 keeps a multi-key citation from printing every author;
%  maxbibnames=99 with giveninits stops the reference list collapsing long
%  author lists to "et al." in a way that would not match the full lists
%  written out in shared.bib (see emmert2021nrlmsis, hersbach2020era5).
%
%  The per-entry reading notes in the bibliography live in the biblatex
%  `annotation' field. They are the most valuable content in that file and are
%  deliberately NOT printed. Note that `annotation' is a DATA field, not an
%  option of \printbibliography -- passing annotation=false to \printbibliography
%  is an xkeyval error, not a way of suppressing it. Blanking the field's format
%  is the mechanism, and it is the right one: the notes survive in the .bib and
%  in the .bbl for the author's use, and never reach the page. This must come
%  AFTER the \usepackage line below, since \DeclareFieldFormat is defined there.
\usepackage[backend=biber,style=numeric,sorting=none,sortcites,
  maxbibnames=99,giveninits=true,maxcitenames=2]{biblatex}

% biblatex requires csquotes whenever babel is loaded, and warns without it
% ("'babel/polyglossia' detected but 'csquotes' missing"). It supplies the
% quote punctuation the bibliography macros expect. Not currently used by any
% entry, but the dependency is unconditional.
\usepackage{csquotes}

\DeclareFieldFormat{annotation}{}

% The thesis bibliography. \addbibresource must be issued in the preamble, and
% this preamble is \input from a directory one level below manuscript/, so the
% path is written relative to that directory, not to this file.
\addbibresource{../thesis.bib}

\usepackage[nameinlink,capitalise]{cleveref}

% ---- Running header: each paper calls \runninghead{Short Title} ----
\newcommand{\runninghead}[1]{%
  \pagestyle{fancy}%
  \fancyhf{}%
  \lhead{\small #1}%
  \rhead{\small\thepage}%
  \renewcommand{\headrulewidth}{0.3pt}%
}
) to compile a submission
%  copy in which all three vanish, exactly as [disable]{todonotes} removes
%  the \todo list --- the boxes must not survive into a submitted draft,
%  and a single switch is the only reliable way to be sure of that.
\newif\ifdraftnotes
\draftnotestrue

\newtcolorbox{decide@box}[1][]{colback=orange!4,colframe=orange!45!black,
  boxrule=0.6pt,arc=2pt,left=6pt,right=6pt,top=4pt,bottom=4pt,#1}
\newtcolorbox{gap@box}[1][]{colback=red!4,colframe=red!50!black,
  boxrule=0.6pt,arc=2pt,left=6pt,right=6pt,top=4pt,bottom=4pt,#1}
\newtcolorbox{idea@box}[1][]{colback=green!4,colframe=green!35!black,
  boxrule=0.6pt,arc=2pt,left=6pt,right=6pt,top=4pt,bottom=4pt,#1}

\ifdraftnotes
  \newenvironment{decide}{\begin{decide@box}}{\end{decide@box}}
  \newenvironment{gap}{\begin{gap@box}}{\end{gap@box}}
  \newenvironment{idea}{\begin{idea@box}}{\end{idea@box}}
\else
  \excludecomment{decide}
  \excludecomment{gap}
  \excludecomment{idea}
\fi

% ---- Theorem-environment aliases ----
%  The drafts write \begin{assumption} and \begin{lemma} in full; the
%  counters above are declared as assum and lem. Alias rather than rename,
%  so neither spelling is a compile error and no draft has to be rewritten
%  to match a decision about abbreviation.
\makeatletter
\@ifundefined{assumption}{%
  \newenvironment{assumption}[1][]{\begin{assum}[#1]}{\end{assum}}}{}
\@ifundefined{lemma}{%
  \newenvironment{lemma}[1][]{\begin{lem}[#1]}{\end{lem}}}{}
\makeatother

\usepackage{babel}

% ---- Cross-references and links (load late: hyperref then cleveref) ----
\usepackage[hidelinks,colorlinks=true,
  linkcolor=blue!55!black,citecolor=blue!55!black,urlcolor=blue!55!black]{hyperref}

% ---- Bibliography ----
%  biblatex, not the plain BibTeX \bibliography. Loaded after hyperref and
%  before cleveref, which is the order the biblatex manual asks for.
%
%  style=numeric keeps the bare-number in-text appearance the papers have
%  always used. sorting=none orders the reference list by first citation,
%  which reads better than alphabetical in a proposal.
%
%  sortcites is the one non-obvious option. It sorts the keys WITHIN each
%  citation so the numeric style can collapse consecutive labels into a
%  range -- without it a five-key \parencite renders as five separate
%  brackets. Note this is NOT natbib's sort&compress, which biblatex does
%  not provide: the biblatex natbib compatibility layer (blx-natbib.def)
%  defines \citep and \citet but nothing for compression, and an unknown
%  package key only warns, so a wrong setting here fails silently.
%
%  maxcitenames=2 keeps a multi-key citation from printing every author;
%  maxbibnames=99 with giveninits stops the reference list collapsing long
%  author lists to "et al." in a way that would not match the full lists
%  written out in shared.bib (see emmert2021nrlmsis, hersbach2020era5).
%
%  The per-entry reading notes in the bibliography live in the biblatex
%  `annotation' field. They are the most valuable content in that file and are
%  deliberately NOT printed. Note that `annotation' is a DATA field, not an
%  option of \printbibliography -- passing annotation=false to \printbibliography
%  is an xkeyval error, not a way of suppressing it. Blanking the field's format
%  is the mechanism, and it is the right one: the notes survive in the .bib and
%  in the .bbl for the author's use, and never reach the page. This must come
%  AFTER the \usepackage line below, since \DeclareFieldFormat is defined there.
\usepackage[backend=biber,style=numeric,sorting=none,sortcites,
  maxbibnames=99,giveninits=true,maxcitenames=2]{biblatex}

% biblatex requires csquotes whenever babel is loaded, and warns without it
% ("'babel/polyglossia' detected but 'csquotes' missing"). It supplies the
% quote punctuation the bibliography macros expect. Not currently used by any
% entry, but the dependency is unconditional.
\usepackage{csquotes}

\DeclareFieldFormat{annotation}{}

% The thesis bibliography. \addbibresource must be issued in the preamble, and
% this preamble is \input from a directory one level below manuscript/, so the
% path is written relative to that directory, not to this file.
\addbibresource{../thesis.bib}

\usepackage[nameinlink,capitalise]{cleveref}

% ---- Running header: each paper calls \runninghead{Short Title} ----
\newcommand{\runninghead}[1]{%
  \pagestyle{fancy}%
  \fancyhf{}%
  \lhead{\small #1}%
  \rhead{\small\thepage}%
  \renewcommand{\headrulewidth}{0.3pt}%
}
) to compile a submission
%  copy in which all three vanish, exactly as [disable]{todonotes} removes
%  the \todo list --- the boxes must not survive into a submitted draft,
%  and a single switch is the only reliable way to be sure of that.
\newif\ifdraftnotes
\draftnotestrue

\newtcolorbox{decide@box}[1][]{colback=orange!4,colframe=orange!45!black,
  boxrule=0.6pt,arc=2pt,left=6pt,right=6pt,top=4pt,bottom=4pt,#1}
\newtcolorbox{gap@box}[1][]{colback=red!4,colframe=red!50!black,
  boxrule=0.6pt,arc=2pt,left=6pt,right=6pt,top=4pt,bottom=4pt,#1}
\newtcolorbox{idea@box}[1][]{colback=green!4,colframe=green!35!black,
  boxrule=0.6pt,arc=2pt,left=6pt,right=6pt,top=4pt,bottom=4pt,#1}

\ifdraftnotes
  \newenvironment{decide}{\begin{decide@box}}{\end{decide@box}}
  \newenvironment{gap}{\begin{gap@box}}{\end{gap@box}}
  \newenvironment{idea}{\begin{idea@box}}{\end{idea@box}}
\else
  \excludecomment{decide}
  \excludecomment{gap}
  \excludecomment{idea}
\fi

% ---- Theorem-environment aliases ----
%  The drafts write \begin{assumption} and \begin{lemma} in full; the
%  counters above are declared as assum and lem. Alias rather than rename,
%  so neither spelling is a compile error and no draft has to be rewritten
%  to match a decision about abbreviation.
\makeatletter
\@ifundefined{assumption}{%
  \newenvironment{assumption}[1][]{\begin{assum}[#1]}{\end{assum}}}{}
\@ifundefined{lemma}{%
  \newenvironment{lemma}[1][]{\begin{lem}[#1]}{\end{lem}}}{}
\makeatother

\usepackage{babel}

% ---- Cross-references and links (load late: hyperref then cleveref) ----
\usepackage[hidelinks,colorlinks=true,
  linkcolor=blue!55!black,citecolor=blue!55!black,urlcolor=blue!55!black]{hyperref}

% ---- Bibliography ----
%  biblatex, not the plain BibTeX \bibliography. Loaded after hyperref and
%  before cleveref, which is the order the biblatex manual asks for.
%
%  style=numeric keeps the bare-number in-text appearance the papers have
%  always used. sorting=none orders the reference list by first citation,
%  which reads better than alphabetical in a proposal.
%
%  sortcites is the one non-obvious option. It sorts the keys WITHIN each
%  citation so the numeric style can collapse consecutive labels into a
%  range -- without it a five-key \parencite renders as five separate
%  brackets. Note this is NOT natbib's sort&compress, which biblatex does
%  not provide: the biblatex natbib compatibility layer (blx-natbib.def)
%  defines \citep and \citet but nothing for compression, and an unknown
%  package key only warns, so a wrong setting here fails silently.
%
%  maxcitenames=2 keeps a multi-key citation from printing every author;
%  maxbibnames=99 with giveninits stops the reference list collapsing long
%  author lists to "et al." in a way that would not match the full lists
%  written out in shared.bib (see emmert2021nrlmsis, hersbach2020era5).
%
%  The per-entry reading notes in the bibliography live in the biblatex
%  `annotation' field. They are the most valuable content in that file and are
%  deliberately NOT printed. Note that `annotation' is a DATA field, not an
%  option of \printbibliography -- passing annotation=false to \printbibliography
%  is an xkeyval error, not a way of suppressing it. Blanking the field's format
%  is the mechanism, and it is the right one: the notes survive in the .bib and
%  in the .bbl for the author's use, and never reach the page. This must come
%  AFTER the \usepackage line below, since \DeclareFieldFormat is defined there.
\usepackage[backend=biber,style=numeric,sorting=none,sortcites,
  maxbibnames=99,giveninits=true,maxcitenames=2]{biblatex}

% biblatex requires csquotes whenever babel is loaded, and warns without it
% ("'babel/polyglossia' detected but 'csquotes' missing"). It supplies the
% quote punctuation the bibliography macros expect. Not currently used by any
% entry, but the dependency is unconditional.
\usepackage{csquotes}

\DeclareFieldFormat{annotation}{}

% The thesis bibliography. \addbibresource must be issued in the preamble, and
% this preamble is \input from a directory one level below manuscript/, so the
% path is written relative to that directory, not to this file.
\addbibresource{../thesis.bib}

\usepackage[nameinlink,capitalise]{cleveref}

% ---- Running header: each paper calls \runninghead{Short Title} ----
\newcommand{\runninghead}[1]{%
  \pagestyle{fancy}%
  \fancyhf{}%
  \lhead{\small #1}%
  \rhead{\small\thepage}%
  \renewcommand{\headrulewidth}{0.3pt}%
}
